\documentclass[12pt]{article}

%% Codierung
\usepackage[utf8]{inputenc}
%% Sprache
\usepackage[ngerman]{babel}
%% Anpassung der Seitenränder
\usepackage[a4paper,lmargin={2cm},rmargin={2cm},tmargin={2.0cm},bmargin = {2.5cm}]{geometry}
%% zum einbinden von Bildern
\usepackage{graphicx} 

\newcommand\circledmark[1][green]{%
  \ooalign{%
    \hidewidth
    \kern0.65ex\raisebox{-0.9ex}{\scalebox{3}{\textcolor{#1}{\textbullet}}}
    \hidewidth\cr
    $\checkmark$\cr
  }%
}

\newcommand\TODO[1]{\textcolor{red}{#1}}

\usepackage{float}
%% und zum positionieren der Bilder
\usepackage[export]{adjustbox}
%% Tabellen
\usepackage{array}
%% Einfärbung von Text
\usepackage{color} 
\usepackage{enumitem}
%% findet Verwendung im Inhaltsverzeichnis
\usepackage[hyphens]{url}
%% versteckt die Hyperref-Links als normalen Text
%% Informationen werden in PDF-Eigenschaften geschrieben
\usepackage[%
hidelinks,%
  pdftitle={Entwicklung eines Physikspiels, welches Physik spielerisch lehrt, auf Basis einer Engine},%
  pdfauthor={Roman Jung und Matthias Seyfarth},%
  pdfsubject={Entwicklung eines Physikspiels},%
  pdfproducer={PDFLaTeX},%
  pdfkeywords={Game; Spiel; lernen; lehren; Counter Strike; Engine; Unity; Studienarbeit; DHBW; Karlsruhe; Informatik}]{hyperref}
%% ermöglicht Querformat
\usepackage{pdflscape}
%% Quellcodelistings
\usepackage{listings}
\lstset{basicstyle=\footnotesize\ttfamily,breaklines=false} %monospace in listings
%Weil latex Kacke ist
\lstset{literate=%
  {Ö}{{\"O}}1
  {Ä}{{\"A}}1
  {Ü}{{\"U}}1
  {ß}{{\ss}}2
  {ü}{{\"u}}1
  {ä}{{\"a}}1
  {ö}{{\"o}}1
}
%für Suche von Umlauten
\usepackage[T1]{fontenc} 
%% ermöglicht den Haken hinter den TODOs
\usepackage{amssymb}
%% für das Abkürzungsverzeichnis
\usepackage{acronym}

\begin{document}
\sloppy
\begin{titlepage}
\includegraphics[height=2.35cm, left]{res/dhbw-logo.jpg}
	\begin{center}
	\linespread{1.5}
	\vspace{0.5cm}
		{\Large\bfseries{Studienarbeit} \par}
		\vspace{1.5em}
		{\large\bfseries {Entwicklung eines Physikspiels, welches Physik spielerisch lehrt,\\ auf Basis einer Engine} \par}

		\vspace{2.5em}
		
		\normalsize{
			für die Prüfung zum\\
			\vspace{0.5em}
			Bachelor of Science\\
			\vspace{0.5em}
			des Studiengangs Informatik\\
			Studienrichtung Angewandte Informatik\\
			\vspace{0.5em}
			an der\\
			\vspace{0.5em}
			Dualen Hochschule Baden-Württemberg Karlsruhe\\
			\vspace{0.5em}
			von\\
			\vspace{0.5em}
			Roman Jung \\
			Matthias Seyfarth\\
			\vspace{7em}
			
			\begin{tabular}{ll}
					Kurs: & TINF13B1 \\
					Dozent: & Prof. Kay Berkling PhD \\
					Abgabedatum: &  XX.XX.2016\\
			\end{tabular}
		}
		\linespread{1.0}
	\end{center}
%\thispagestyle{empty}    
\end{titlepage}
\setcounter{page}{2}
\newpage

\begin{abstract}
Diese Studienarbeit befasst sich mit der Entwicklung eines Puzzle-Spiels, welches Physik spielerisch lehrt. Hierbei werden wissenschaftliche Studien analysiert, die das Spiel als Lernmedium untersuchten. Basierend auf diesen Analysen wird ein Spielkonzept erstellt und programmiert. 

Einerseits wird sich mit der Thematik befasst, ob das Spiel in 2D oder 3D erstellt wird. Entsprechend dazu werden unterschiedliche Spiele-Engines betrachtet und analysiert, um die passendste zu ermitteln. 

Nach der Wahl / Vor der Wahl einer Engine wird sich mit den unterschiedlichen Bereichen der Physik befasst und eines als Grundlage des Spiels gewählt (Wir würden Unity als Engine verwenden).\\

Ein Spielkonzept könnte folgendermaßen aussehen:
\begin{itemize}
\item{Thematik:}  Licht
\item{Grundidee:} Ein Lichtstrahl soll durch Bewegungen auf ein Ziel gelenkt werden.
\item{Ziel:} Das Licht auf eine oder mehrere Knospen lenken, die zu einer Blume (o.ä.) heranwachsen und das Level abschließen
\item{optional:} Sammelobjekte vor dem Abschluss des Levels einsammeln, eine Punktesystem
\item{Einflüsse:} Manuell setzbare Spiegel (Reflexion), Flüssigkeiten(Brechung) und weitere Objekte aus der Lichtphysik 
\item{Features:} Je nach Anzahl der gesetzten Objekte und den gesammelten Objekten entsteht ein Ranking

    \begin{itemize}
        \item Schnittstelle: Erweiterbar mit weiteren physikalischen Bereiche
        \item Schnittstelle: die einzelnen Bereiche (z.B. Licht, Schwerkraft, Magnetismus) mit neuen Lv erweitern können
    \end{itemize}

\end{itemize}

\end{abstract}
\newpage

\thispagestyle{empty}
\tableofcontents
\thispagestyle{empty}
\newpage

% Abbildungsverzeichnis
\addcontentsline{toc}{section}{Abbildungsverzeichnis} 
\listoffigures 
% Tabellenverzeichnis
\addcontentsline{toc}{section}{Tabellenverzeichnis}
\listoftables
\newpage

\phantomsection \addcontentsline{toc}{section}{Abkürzungsverzeichnis}
\renewcommand\refname{Abkürzungsverzeichnis} \section*{Abkürzungsverzeichnis}
\begin{acronym}[XXXXXXXXX]
 \acro{UE4}{Unreal Engine 4}
 \acro{U5}{Unity 5}
 \acro{BGE}{Blender Game Engine}

\end{acronym}

\newpage
\newpage

%% input Kapitel
\section{Einleitung}

\subsection{Themenabgrenzung und Zielsetzung}
\subsection{Methodik und Aufbau der Arbeit}



\newpage
\section{Analyse}
\subsection{Was ist ein Spiel?}
\subsection{Lernen durch Videospiele}
Computerspiele leiden unter vielen schlechten Vorurteilen. Ihnen wird eine Steigerung der Aggressivität, eine Vereinsamung oder das Verdummen der Spieler angelastet. Studien widerlegen diese Vorurteile vehement \cite{psylex.de2012}. Entgegengesetzt dazu weisen einige Videospiele sogar Lerneffekte auf \cite{paradisiredaktion2010}. Doch was genau lehren sie \TODO{Warum/Wieso lehren Sie, Wie funktioniert es}? Um dieser Frage nach zu kommen werden einige Spiele genauer betrachtet, dessen Inhalte analysiert und Lernkompetenzen aufgezeigt.

\subsubsection[Age of Empires]{Age of Empires\footnote{Dieser Abschnitt wird an mehreren Stellen von der Quelle \cite{bundeszentralefuerpolitischebildung2005} referenziert.}}
Eines der betrachteten Spiele ist "{}Age of Empires"{}. Age of Empires ist ein Echtzeit-Strategiespiel, welches sehr gute Grafik und Animationen aufweist. Das Ziel ist es, alle gegnerischen Völker auszulöschen und dabei das Überleben des eigenen Volkes sicherzustellen. Um dieses Ziel zu erreichen, muss ein Imperium erschaffen werden, welches eine ausgeklügelte Infrastruktur und Militärwesen besitzt. Dazu wird eine Vielzahl unterschiedlicher Rohstoffe benötigt, die für die Erstellung der verschiedenen Gebäude, Schiffe und Soldatentypen notwendig sind. Das Ressourcenmanagement ist daher ein essentieller Bestandteil des Spielgeschehens.

Dem Spieler stehen verschiedenste Handlungs- und Entscheidungsmöglichkeiten zur Verfügung. Diese fordern ihn in taktischer und strategischer Disziplin heraus. Aus diesem Grund benötigt dieser eine komplexe denkerische Fähigkeit, um abschätzen zu können, welche Schritte für das eigene Überleben und die Auslöschung des Feindes notwendig sind. Daneben muss er in der Lage sein die Spielsituation zu erfassen und die Strategie des Gegners in die eigene Vorgehensweise einzuberechnen. Zuzüglich sind Stressbewältigung und Reaktionsgeschwindigkeit von großer Bedeutung, der auf dem "{}Realtime-Modus"{} zurückzuführen ist. 

Die meisten Spiele vom Genre Echtzeit-Strategie weisen eine Levelstruktur auf. Mit zunehmenden Levelfortschritts werden die Anforderungen an den Spieler vielfältiger und komplexer. Aufgrund dessen wird der Denkprozess des Spielers angeregt. Neue Methoden und Strategien müssen ausgebildet werden, um das Ziel zu erreichen. Spieler, die bereits Erfahrungen mit diesem Genre gesammelt haben, konnten nachweislich besser Probleme lösen als unerfahrene Spieler. Außerdem waren sie besser in der Lage neue Probleme zu erkennen und diese zu verarbeiten. Die erfahrenere Gruppe der Spieler erkannten Zusammenhänge zwischen den einzelnen Spielelementen schneller und legten Prioritäten, zum Beispiel auf den Ausbau der Stadt. Die Unerfahrenden testeten die unterschiedlichen Elemente aus. Erst beim wiederholten Versuch des selben Levels konnte es erfolgreich abgeschlossen werden. Das wurde durch das Erkennen der Zusammenhänge der unterschiedlichen Spielelemente erreicht.

Der Spieler muss in Age of Empires unterschiedliche Probleme lösen. Darunter zählt ein ausgeklügeltes Ressourcenmanagement, um die verschiedenen Gebäude- und Soldatentypen ausbilden zu können. Darüber hinaus müssen im fortgeschrittenen Levelabschnitt neue Taktiken und Handlungsschemata ausgebildet werden, um die komplexer werdenden Aufgaben zu erfüllen und den Feind auszulöschen. Darüber hinaus wird die allgemeine Wahrnehmung geschult. Das zeigt sich durch frühzeitige Erkennung der feindlichen Aktivitäten, woraus sich wiederum eine erhöhte Reaktionsgeschwindigkeit ableitet, die zum Einleiten der Gegenmaßnahmen notwendig sind. 

Ein Lerneffekt ist nach Zusammenfassung dieser Aspekte in Age of Empires nachgewiesen.

\subsubsection{Counter Strike}
Ein weiteres untersuchtes Spiel ist "{}Counter Strike"{}. Hierbei handelt es sich um ein 3D Online-Taktik-Shooter, wo eine Terrorristengruppen gegen eine Antiterroreinheit kämpft. Das Ziel ist es, das jeweils gegnerische Team vollständig auszulöschen oder die gesetzte Bombe der Terrorristen zu entschärfen bzw. die Geiseln zu retten. Zur Verfügung steht ein großes Repertoire an Gewehren, Pistolen, Messer, aber auch Schutzkleidung, welches mit Geld erworben wird.

Um das Ziel des Spiels zu erreichen muss sich der Spieler bewegen können. Das ist allerdings nur möglich, wenn er sich in einem 3-Dimensionalem Raum zurechtfinden kann. Dafür wird die räumliche Vorstellungskraft benötigt, die sich während des Spielens verbessert. Objekte und Hindernisse werden mit zunehmender Spielerfahrung leichter überwunden, auch die Schätzungen von Entfernungen sind präziser. Vor allem die kognitive Wahrnehmung der Umgebung, insbesondere der bewegte Objekten werden trainiert. 

Neben der Bewegung ist auch das Zielen ein ausschlaggebender Bestandteil des Spielgeschehens. Die Maus, welche direkt mit der Zielvorrichtung gekoppelt ist, muss zum Gegner gesteuert werden. Dieses Verhalten fordert eine ausgesprochen hohe Hand-Augen-Koordination. In Abhängigkeit der Lokation des Gegners (Das Auge nimmt den Gegner als Reiz auf.) muss die Maus in Richtung des Ziels gesteuert werden (Die Hand steuert die Maus zur gewünschten Stelle, sodass sich der Feind im Bildzentrum befindet.). Neben der Hand-Augen-Koordination ist auch dort die räumliche Vorstellungskraft wesentlicher Bestandteil, weil sie das Anvisieren erst ermöglicht. 

Feinde können sich blitzschnell um Ecken bewegen und den Spieler erschießen. Um das zu verhindern kann sich der Spieler hinter Objekte verstecken oder schnellstmöglich den Feind auslöschen. Das ist allerdings nur möglich, wenn eine hohe Reaktionsgeschwindigkeit besteht. Unerfahrende Spieler weisen in der Regel ein langsameres Reaktionsverhalten auf und werden häufiger von der gegnerischen Fraktion getötet. Erfahrene Spieler, welche bereits viel Spielzeit aufgebracht haben, zeigen ein deutlich erhöhtes Reaktionsverhalten und können innerhalb weniger Augenblicke Entscheidungen treffen, die ihr Überleben sichern.

Entscheidende Faktoren, die zur Verbesserung der Reaktion beitragen, sind geschulte visuelle und akustische Aufmerksamkeit. Die visuelle Aufmerksamkeit wird hauptsächlich durch die Bewegung andere Spieler verursacht. Das liegt an der evolutionär bedingten Wahrnehmung des Auges, welches sofortige Aufmerksamkeit zum bewegten Objekt verursacht \cite{juliagross2011}. Neben der visuellen kann ebenfalls die akustische Wahrnehmung geschult werden. Personen, die dieses Spiel mit Ton spielen, haben die Möglichkeit andere Charaktere zu hören. Dabei sind Explosionen, Schüsse, aber auch die schnellen Bewegungen (Rennen des Charakters) eines Spielers wahrzunehmen. Diese Erkenntnis kann genutzt werden, um den Feind zu überraschen oder um Taktiken auszuarbeiten.

Counter Strike ist ein Multiplayer-Spiel. Das bedeutet, es treffen mehrere reale Personen aufeinander, die ein Team bilden. Ein erfolgreiches Abschließen der Mission kann durch Teamarbeit und koordiniertes Agieren der Spieler vereinfacht werden. Das erfordert Kommunikationen zwischen den Teammitgliedern, welche in einem Chat oder über einem Sprachdienst abgehalten wird. Spieler, die oft einen der Dienste in Anspruch nehmen, schulen somit ihre kommunikative Fähigkeiten und verbessern unter Umständen ihre Fremdsprachenkenntnisse.

Der Spieler muss in Counter Strike die gegnerische Fraktion vollständig auslöschen oder die Nebenmission erfüllen. Um diese Ziele zu erreichen, müssen unterschiedliche Fähigkeiten und Kompetenzen genutzt werden. Einerseits ist eine umfangreiche Hand-Augen-Koordination notwendig, damit der Feind anvisiert werden kann. Darüber hinaus ist eine ausgebildete räumliche Vorstellungskraft unverzichtbar, welche für die Bewegung, aber auch das Zielen ermöglicht. Ein trainiertes Reaktionsverhalten ist ebenfalls essentiell. Es dient zum eigenen Schutz, aber auch um blitzschnell zu reagieren, um den Feind auszulöschen. Unterstützt wird diese Fähigkeit durch eine gute visuelle und akustische Wahrnehmung, welche sich durch zunehmender Spielerfahrung sensibilisiert.

Nach Zusammenfassung dieser Aspekte ist auch in Counter Strike ein Lernprozess nachweisbar.

\subsubsection{Flappy Bird}
Das zuletzt untersuche Spiel ist ein Mobile Game, welches zum Genre Jump $'$n Run zählt. Hierbei handelt es sich um das bekannte Spiel "{}Flappy Bird"{}. Flappy Bird ist ein 2D-Spiel, welches eine ähnliche Grafik wie Super Mario World aufweist. Das Ziel des Spiels ist es, einen Vogel zwischen zwei Gegenstände zu führen und möglichst viele Punkte zu sammeln. Die Welt bewegt sich dabei von rechts nach links.

Um das Spiel zu starten muss der Bildschirm berührt werden, worauf hin der Vogel nach unten fällt. Durch erneutes berühren des Touchscreens "springt" der Vogel nach oben, um kurze Zeit später wieder nach unten zu fallen. Um diese Aufgabe zu erfüllen ist eine große Hand-Augen-Koordination notwendig, da sich der Vogel im richtigen Zeitpunkt nach oben bewegen muss. Berührt dieser eines der Hindernisse so ist das Spiel beendet und es kann erneut begonnen werden.

Das Mobile Spiel Flappy Bird weist ebenfalls einen lehrenden Aspekt auf. Allerdings beschränkt sich die Ausbildung nur auf die Hand-Augen-Koordination, da sie der einzige und wichtigste Aspekt darstellt.

\subsubsection{Lernen durch Videospiele: Fazit}
"{}Age of Empires"{}, "{}Counter Strike"{} und "{}Flappy Bird"{} sind Spiele, die aus unterschiedliche Genres stammen. Sie stellen verschiedene Anforderungen an den Spieler, verwenden differenzierte Spielprinzipien und weisen verschiedenartige Ziele auf. Doch sie besitzen eine Gemeinsamkeit: den Spieler zu motivieren die Feinde zu vernichten, möglichst viele Punkte oder Erfolge zu sammeln. Diese Motivation veranlasst den Spieler das Spiel erneut zu spielen, sodass dieser mehr Punkte sammelt und sich selbst übertrifft. Direktes Feedback vom Spiel, z.B. das Erhalten von Punkte in Flappy Bird, spornt an, noch mehr Punkte zu erreichen. 

Durch das Wiederholen der Spieltätigkeit mit dem Ziel besser zu sein als vorher oder den Gegner bezwingen zu können, werden die einzelnen Kompetenzen, wie analytisches Denken, Reaktionsgeschwindigkeit, visuelle und akustische Wahrnehmung geschult und stetig verbessert. Auch das Zurechtfinden in einem 3D-Raum, die Hand-Augen-Koordination oder soziale Qualitäten, wie Teamfähigkeit oder Kommunikation in einer Gruppe werden gefördert und hochwertiger. 

Obwohl alle der oben genannten Spiele verschiedene Genres bedienen und Tätigkeiten abverlangen, haben diese doch eines gemeinsam: sie vermitteln alle eine spezielle Art von Wissen, die sich direkt auf unsere Fähigkeiten auswirken. 

\subsection{Wahl der Spielart}
\TODO{Hier kommt Text warum wir uns für das Lichtspiel entschieden haben}
\TODO{Wie sind wir auf die Idee gekommen?}
\TODO{Idee wurde unterstützt durch das Spiel God of Lights}

\subsection{Spielgenre}
\TODO{Hier kommt Text warum wir uns für das Genre entschieden haben}

\subsection{Spielkonzept}
\subsubsection{Grundidee}
Die Grundidee dieses Spiels ist es, einen Weltenbaum mit Licht zu versorgen und dadurch wachsen zu lassen. Dafür muss der Lichtstrahl durch unterschiedliche Puzzle-Rätsel ans Ziel geleitet werden. Um das Ziel zu erreichen muss mit unterschiedlichen Objekten interagiert werden. Mit zunehmenden Levelfortschritts müssen komplexere Objekte verwendet werden. Im Laufe des Spiels entwickelt sich der Samen zum Weltenbaum.

\subsubsection{Darstellung des Spiels}
Das Spiel soll in einem Comichaften 2.5D Stil dargestellt werden. Dabei sollen die Grafiken in einem tristen Schwarzweiß mit farbigen Hintergrund verwendet werden. Nach dem Erleuchten des einzelnen Gebiets sollen diese in bunten Farben erstrahlen. Die einzelnen Szenarien und Level befinden sich in Höhlen, auf Wiesen und im Wald.

\TODO {HIER GRAFIKELEMENTE AUFZEIGEN}


Die Idee und Darstellung des Spiels wird im Kapitel \ref{ganmeDesignDoku} genauer erläutert.

\subsection{Engine}
\subsubsection{Was ist eine Engine?}
Eine Spielengine ist der „Motor“ hinter dem eigentlichen Spiel und wird als Framework bzw. Entwicklungsumgebung angeboten. Die Engine steuert den gesamten Spielverlauf und ist ebenfalls für die visuelle Darstellung sowie die Berechnung der Physik verantwortlich. In folgenden Kapitel werden mehrere unterschiedliche Engines verglichen und anhand unserer Anforderungen eine passende ausgewählt. 

\subsubsection{\ac{UE4} vs \ac{U5} vs \ac{BGE}}[Test]
Da es sich bei dem Spiel um ein Puzzle und Logikspiel handelt, mit welchem auch Lerninhalte vermittelt werden sollen, werden die Casualgamer und Kinder als Zielgruppe ausgewählt. Um diese Zielgruppe zu erreichen muss die entsprechende Plattform gewählt werden, welche eine möglichst große Anzahl der Zielgruppe abdeckt. Dafür eignen sich besonders Mobilgeräte, da diese heutzutage fast jeder besitzt und auch darauf Spiele spielt.  \\

\textbf{Plattformen:}\\
Alle zu vergleichenden Engines unterstützen Windows und Linux Rechner sowie MAC OS X. Zusätzlich sind \ac{U5} und \ac{UE4} mit den gängigen mobilen Betriebssystemen Android und Apples IOS kompatibel und erlauben auch die Entwicklung an Microsofts und Sonys aktuellen Konsolen Xbox One und Playstation 4. Doch auch ganz neue Betriebssysteme und Technologien wie VR und Steam OS werden von Haus aus unterstützt.

\begin{table}[H]
\centering
\renewcommand{\arraystretch}{1.5}
\begin{tabular}{lp{14.5cm}}
\ac{U5}:  & Zusätzlich unterstützt \ac{U5} weitere mobile Plattformen wie Windows Phone 8, Tizen und die Playstation Vita. Auch für TV Geräte wie Android TV oder Samsung Smart TV können Spiele mit der \ac{U5} Engine entwickelt und gespielt werden. Im Gegensatz zu \ac{UE4} ist \ac{U5} auch mit älterer Konsolenhardware, wie der Xbox 360 und der Wii U kompatibel. Aber auch Webanwendung mit WebGL oder dem Unity Web Player können erstellt werden. \\

\ac{UE4}: & \ac{UE4} hat den weiteren Vorteil dass auch HTML5 unterstützt wird. \\

\ac{BGE}:     & Mit \ac{BGE} ist auch die Entwicklung für Mobilgeräte mit Android Betriebssystem möglich. Voraussetzung dafür ist, dass auf dem Mobilgerät Android \ac{BGE} Player GSOC 2012 installiert ist. 
\end{tabular}
\end{table}

Wie man sehen kann erfüllen alle Engines die Voraussetzung mit mobilen Geräten kompatibel zu sein. Da es sich bei diesem Projekt um ein Studienprojekt handelt, sollten sich auch die Lizenzgebühren und die laufenden Kosten im unteren Bereich befinden. Womit der nächste Kritikpunkt unseres Vergleiches die Lizenzierung sowie die Kosten betreffen.


\newpage

\textbf{Lizenz und Kosten:}
\begin{table}[H]
\centering
\renewcommand{\arraystretch}{1.5}
\begin{tabular}{lp{14.5cm}}
\ac{U5}: & \ac{U5} kann als zwei verschiedenen Arten erworben werden. Die Personal Edition steht dabei Kostenfrei zur Verfügung und kann mit vollen Featureumfang der Engine erworben werden. Zusätzlich gibt es die Pro Variante, welche pro Release gezahlt oder durch einen Monatlichen Beitrag erworben werden kann. Diese Edition bietet weitere Möglichkeiten effizienter zu arbeiten, wie zum Beispiel durch eine Teamlizenz. Weiterhin gibt es in der Pro Variante mehr Möglichkeiten sein erstelltes Spiel durch eigene Bootanimationen u.ä. zu gestalten.\\
UE4:& Die Unreal Engine ist seit März 2015 kostenfrei zu erwerben und kann bedingungslos verwendet werden. Falls beim Verkauf des entwickelten Spiels mehr als 3000 Dollar Umsatz erwirtschaftet wird, hält sich Epic Games (die Entwickler der Unreal Engine) das Recht vor, 5 Prozent des Umsatzes als Lizenzgebühr zu  verlangen.\\
\ac{BGE}: & \ac{BGE} ist eine freie Software  und kann uneingeschränkt genutzt, verändert und die erstellten Produkte verkauft werden.
\end{tabular}
\end{table}
Alle der genannten Engines sind für dieses Projekt ohne Kosten und Lizenzgebühren nutzbar und eignen sich für diese Arbeit.
Für dieses Projekt ist nicht nur eine Kostenübersicht notwendig, sondern sollte auch ein guter Einstieg durch Tutorials/ Literatur/Dokumentationen ermöglicht werden. \\

\textbf{Erlernbarkeit:}\\
Alle der besagten Engines besitzen eine sehr gute Dokumentationen. Durch die visuelle Unterstützung durch Screenshots von der \ac{UE4}  wird die ohnehin gut beschriebene Dokumentation noch hochwertiger. Jede der Engines besitzt eine große Community, wodurch auch sehr viele Tutorials und Wiki Einträge vorhanden sind. Zusätzlich gibt es eine  vielzahl an  Anleitungen von professionellen Entwicklern, die ihr Wissen allen Benutzergruppen (Neuling/Fortgeschritten/Experte) zur Verfügung stellen. Weiterhin gibt es eine große Auswahl an Literatur für den Einstieg mit \ac{U5} und \ac{BGE}. Da sich \ac{UE4} noch in der Entwicklung befindet, ist die Literaturauswahl eingeschränkt (vor allem in der deutschen Sprache).

Alle der Engines haben eine große Community und bieten auch sehr viele gute einsteigerfreundliche Möglichkeiten an. 
Um eine gute Einarbeitung zu ermöglichen und mit der Engine effizient zu arbeiten spielt auch die Nutzbarkeit und die Komplexität des Interface eine Rolle.\\

\textbf{Nutzbarkeit:}\\
Die einzelnen Interfaces unterscheiden sich nur leicht voneinander und sind mit beweglichen/skalierbaren Fenstern flexibel anpassbar. Die einzelnen Entwicklungsumgebungen können auf allen gängigen Betriebssystemen eingesetzt werden.
\begin{table}[H]
\centering
\begin{tabular}{lp{14.5cm}}
\ac{UE4}:& 
Das \ac{UE4} Interface ist komplexer als das von \ac{U5}. Dieses ist eher für Benutzer mit Erfahrung geeignet, da es den Einstieg etwas erschwert. Darüber hinaus beansprucht \ac{UE4} mehr Zeit für das Ausführen von Aufgaben, wie Speichern des Projekts oder Importieren von Assets. Zudem kommt hinzu, dass \ac{UE4} eine sehr hardwarelastige Engine ist, welche mehr Ressourcen des Endgerätes verwendet, als die anderen Engines.\\
\ac{U5}:&
\ac{U5} hat den Vorteil, dass es sehr schnell reagiert und auch kleinere Aufgaben sehr schnell abarbeitet.\\  
\ac{BGE}:&
Das \ac{BGE} Interface ist, wie auch das \ac{UE4} Interface, komplexer und benötigt mehr Einarbeitungszeit.
\end{tabular}
\end{table}
Da das Projekt für Mobilgeräte entwickelt wird, welche nicht immer mit der aktueller Hardware ausgestattet ist, eignet sich gerade aus dem Aspekt der Nutzbarkeit die Unity Engine besonders gut.
Dieses Projekt ist für ein Team dieser Größe nur sehr schwer realisierbar, daher ist die Unterstützung aus der Community wichtig. Die Engines sollten bestimmte Assets zur Verfügung stellen und auch eine möglichst günstige und vielfältige Auswahl anbieten.\\

\textbf{Unterstützung:}\\
Durch sogenannte Assets kann man sich bereits erstellte Projektteile herunterladen und in seiner eigenen Arbeit verwenden. Diese können einfache Scripte bis hin zu komplexe Grafikelemente oder Animationen sein. Hierbei stehen kostenfreie sowie kostenpflichtige zur Verfügung. Da \ac{UE4} relativ neu ist, ist auch die Auswahl kleiner als im Unitystore. Nichtsdestotrotz kommen täglich neue Inhalte hinzu. \ac{BGE} stellt aus ihrer Community eine umfangreiche Auswahl aus allerlei verschiedensten Bereichen von Assets bereit. Diese werden meistens kostenlos angeboten und enthalten oft eine vollständige Anleitung wie dieses Asset erstellt wurde. Bei \ac{U5} sind ebenfalls sehr viele Assets kostenlos. Die kostenpflichtigen fallen in der Regel günstiger aus, als die von \ac{UE4}.

Die Auswahl ist bei allen drei Engines riesig. Jedoch ist durch die längere Etablierung der \ac{U5} und \ac{BGE} eine deutlich größere Auswahlmöglichkeit vorhanden. Diese sind meist günstiger oder sogar komplett Kostenfrei zu erwerben.\\

\textbf{Physik:}\\
Durch Physik entsteht der Eindruck in jeder einzelnen Szene des Spiels interagieren zu können. Der Spieler erhält dadurch ein direktes Feedback. Eine korrekte Physik gibt dem Spieler das Gefühl in die Spielewelt "eintauchen" zu können. Diese kann dadurch als real wahrgenommen werden. Gerade für ein Spiel das mit physikalischen Grundprinzipien arbeitet, ist dies ein wichtiger Bestandteil den man bei der Entscheidung berücksichtigen sollte. 
\begin{table}[H]
\centering
\begin{tabular}{lp{14.5cm}}

\ac{U5}: &
\ac{U5} basiert für physikalische Berechnungen auf der PhysX Engine, welche für 3D Modellierungen verwendet wird. Für 2D Projekte nutzt \ac{U5} die Box 2D Engine. Diese ist Open Source und ist in \ac{U5} bereits nativ implementiert. Box 2D eignet sich speziell für Rigidbody2D Objekte, mit dem auch \ac{U5} arbeitet.\\
\ac{UE4}:& Die Physik in \ac{UE4} basiert ebenfalls auf der PhysX Engine von Nvidia. Diese ermöglicht es, korrektes physikalisches Kollissionsverhalten zwischen mehreren 3D Objekten zu berechnen und darzustellen. Für 2D Projekte verwendet \ac{UE4} die eigen entwickelte Paper2D-Engine. \\

\ac{BGE}:&
Bei \ac{BGE} kommt für die Physik die Bullet Physik Engine zum Einsatz. Diese ermöglicht die korrekte physikalische Berechnungen für Partikelsysteme, starre Körper sowie Flüssigkeiten, Kraftfelder und Gravitation.
\end{tabular}
\end{table}


\textbf{Grafik:}\\
Die Grafik eines Spiels bestärkt die Immersion und sorgt für eine optische Darstellung des Spiels. Dieses Spiel soll im 2.5D stilisiert werden und die Engine sollte entsprechende Grafiken dazu ermöglichen.
\begin{table}[H]
\centering
\begin{tabular}{lp{14.5cm}}

\ac{UE4}:&
Die \ac{UE4} verwendet für das Rendern der Grafiken ein eigenes Rendering System, welches auf OpenGL basiert. \ac{UE4} verwendet eine neue DirectX11-Pipeline, welche es dem Grafikprozessor erlaubt Änderungen von Beleuchtungen und Shading-Effekte im Nachhinein vorzunehmen.\\
\ac{U5}: &
\ac{U5} unterstützt OpenGL und Direct3D. Welches Shadingverfahren dabei zum Einsatz kommt hängt dabei von der  Zielplattform ab. Die integrierten Beleuchtungsmodelle können mit selbsterstellten Shadern erweitert oder modifiziert werden.\\
\ac{BGE}:&
Blender ist eine 3D Software zum Erstellen von Animationen oder Modellen, welche auf DirectX und OpenGL basiert. Zudem ist eine eigen entwickelte Game-Engine integriert. Der Unterschied zwischen dem "normalen Modus" und \ac{BGE} ist der Render-Vorgang. In der Game Engine wird jede Szene einzeln und in Echtzeit gerendert. Im Gegensatz dazu rendert der normale Modus die Modelle/Bilder/Animationen im voraus und sind nach dem Vorgang nicht mehr veränderbar.
\end{tabular}
\end{table}
Dadurch dass die \ac{BGE} nicht ohne weiteres 2D Projekt unterstützt, ist \ac{U5} oder \ac{UE4} die bessere Variante.
Die Entwicklungssprache der Engine ist nur zweitrangig, jedoch sollte diese verglichen werden, um zum Beispiel Einarbeitungszeiten zu verkürzen.\\

\textbf{Scriptsprachen:}
\begin{table}[H]
\centering
\begin{tabular}{lp{14.5cm}}
\ac{UE4}:&
\ac{UE4} setzt bei der Erstellung der Spiele auf C++  und den sogenannten Blueprint Scripting. Mit diesen ist es möglich Abläufe von bestimmten Aktionen /Situationen/Abläufen durch visuelle Darstellung durch Modelle zu definieren und Erstellen.\\
\ac{U5}:&
\ac{U5} verwendet als Scriptsprachen C\#, Boo und Unityscript, welches an Javascript anlehnt und aufbaut. Man ist nicht auf eine Sprache begrenzt und es ist frei, welche man für die Entwicklung verwendet. \ac{U5} bietet die Möglichkeit der visuellen Scripte, wie \ac{UE4} standardmäßig, nicht an. Erweiterungen dafür können aber im Assetstore erworben werden.\\

\ac{BGE}:&
Für die Entwicklung mit der \ac{BGE} wird ein Logikeditor verwendet, welcher Logikbausteine für Interaktionen, Animationen und Situationen anbietet. Diese sind durch C++ oder Python Scripte erweiterbar.

\end{tabular}
\end{table}

\subsubsection{Fazit}
Alle oben genannten Kriterien wurden beachtet, um sich für eine Engine zu entscheiden, die für das eigene Spiel eingesetzt werden soll. Da kaum bis keine Erfahrung mit dem Umgang von Engines besteht, ist ein guter Einstieg und die Möglichkeiten sich durch Tutorials und Literatur einzuarbeiten sehr wichtig. Zusätzlich wurde sich für ein 2.5D Stil der Grafik entschieden, welches mit Lichtphysik interagiert und auf mobilen Plattformen unterstützt wird. Durch die geringe Leistung vieler Geräte ist es wichtig auf die Performance der Geräte zu achten. Diese Voraussetzungen und die Erarbeitung der einzelnen Unterschiede machten schnell klar, dass sich \ac{U5} für dieses Projekt am besten eignet. Es bietet eine umfangreiche Möglichkeit der Einarbeitung, einen 2D Rendermodus sowie Physikeffekte, und unterstützt mobile Endgeräte mit leistungsschwacher Hardware. 

\subsection{Konzepte und Formeln aus der Physik}


\newpage
\section{Game Design Document}
\label{ganmeDesignDoku}
\subsection{Geschichte / Story}
\subsubsection{Ziel des Spiels}
Das Ziel des Spiels ist es, einen Samen mit Licht zu versorgen, sodass dieser zu einem Weltenbaum heranwächst. Dafür werden mehrere Level durchlaufen, in dem ein Lichtstrahl durch unterschiedliche Spielelemente von einem Startpunkt zu einem Zielpunkt (Samen des Baums) geleitet wird.

\subsection{Level und Menüs}
\subsubsection{Menü}
\begin{description}

\item [Hauptmenü] \hfill \\
siehe Zeichnung (nicht eingefügt)

\item [Optionsmenü] \hfill \\
siehe Zeichnung (Play und Optionen Buttons)(nicht eingefügt)

\item [Abschlussdarstellung] \hfill \\
Siehe Zeichnung
Es ist möglich das Level Neuzustarten, das nächste Level zu starten, in die Levelansicht zu wechseln und/oder das gerade gespielte Level komplett zu betrachten. 
\end{description}
\subsubsection{Levelkonzept}
Das Konzept dieses Spiels beruht auf die Verwendung mehrerer aneinandergereihter Bereiche, die den Charakter eines Sidescrollers aufweisen. Zu Beginn des Spiels weisen die Levels nur einen Bereich auf. Diese Level dienen der Orientierung und das Verstehen der unterschiedlichen Spielfunktionen. In den weiterführenden Level steigt die Anzahl der Bereiche, die gelöst werden müssen, um das Ziel zu erreichen. Der Wechsel zwischen den Bereichen erfolgt allerdings nicht autonom. Hierfür muss am Ende eines jeden Bereiches ein Checkpoint mit dem Lichtstahl angeleuchtet werden. Durch zusätzliches berühren dieses Punktes wird der aktuelle Bereich beendet und der neue freigeschaltet. Das zeigt sich durch verschieben der Kamera in eine Richtung, bis sich der Checkpoint am Anfang des Screens befindet. Durch das selbstständige Ausstrahlen des Lichtstrahls fungiert er als eine Art 2. Startpunkt. Die Zugehörigkeiten zwischen die Anzahl der Bereiche und die Level sind in Tabelle \ref{screenLevel} aufgezeigt. Neue Level werden durch den Abschluss des vorherigen Levels und  eine definierte Anzahl an Collectables freigeschaltet. Im Durchschnitt sind 3 Coll. gedacht.

\begin{table}[H]
\centering
\begin{tabular}{c|c}
Levelbereich & Anzahl der Screens / Bereiche \\ \hline
1 - 5  & 1 Screen \\ 
6 - 15 & 2 Screens \\
16 - 25 & 3 Screens      
\end{tabular}
\caption{Anzahl Screens in Levelbereiche}
\label{screenLevel}
\end{table}

\subsubsection{Anzahl Collectables im Level}
Jedes Level besitzt eine genau definierte Anzahl Collectables, welche für die Freischaltung neuer Level benötigt werden. Der Zusammenhang zwischen Collectable, Bereich und Level wird in Tabelle \ref{collect_Level} gezeigt. Insgesamt können somit nach Abschluss des letzten Level 110 Collectables gesammelt werden.
\begin{table}[H]
\centering
\begin{tabular}{l|c|c|c}
Levelbereich & \begin{tabular}[c]{@{}c@{}}Anzahl Collectable \\ pro Bereich\end{tabular} & \begin{tabular}[c]{@{}c@{}}Anzahl Collectable \\ pro Level\end{tabular} & \begin{tabular}[c]{@{}c@{}}Gesamte Collectables \\ dieser Levelspanne\end{tabular}  \\ \hline
Lv 1 - 5 &2  &2  &10   \\ 
Lv 6 - 15&2  &4  &40  \\
Lv 16 - 25 &2  &6  &60
\end{tabular}
\caption{Anzahl Collectables}
\label{collect_Level}
\end{table}

\subsubsection{Setting}
Das Setting beschreibt die Umgebung und die Konzeption der einzelnen Level. Folgende Settings werden verwendet:

\begin{itemize}
\item Höhle
\item Erdoberfläche/Wiese
\item Wald
\item Baumkrone
\item Bauminneres (optional)
\item Himmel (optional)
\end{itemize}

\subsection{Spielelemente}
\begin{enumerate}
\item \textbf{Spiegel} \hfill \\
Ein Spiegel ist ein Game Objekt, welches den Lichtstrahl reflektiert. Hierbei wird die korrekte Reflexion nach physikalischen Gesetzen angewendet. Werden Spiegel von hinten angeleuchtet erfolgt keine Reflexion. Alle Spiegeltypen werden mit einer vordefinierten Richtung im Spiel angebracht.

Der Spiegel ist als abstrakte Form zu sehen und weist deshalb keine grafischen Elemente und Interaktionsmöglichkeiten auf. Diese werden bei den spezifischen Spiegeltypen definiert.

\begin{enumerate}
    \item \textbf{Rotierender Spiegel} \hfill \\
    Der rotierende Spiegel befindet sich an einer definierten Stelle im Spiel. Durch antippen des Game Objektes wird dieser markiert und kann durch wischen um insgesamt 140$\degree$ verdreht werden. Der rückseitige Mittelpunkt befindet sich an der Wand und stellt die Verbindung mit der Wand sicher. Der Lichtstrahl wird entsprechend des Auftrittswinkels auf der Spiegelfäche reflektiert.
    
    \TODO{TODO: Aussehen des Spiegels}
    
    \item \textbf{Spiegel mit Farbeinschränkungen} \hfill \\
    Der Spiegel mit Farbeinschränkungen kann als fester, rotierender oder beweglicher Spiegel auftreten. Die Funktionalität sowie die Interaktionsmöglichkeiten beschränken sich auf die Spiegeltypen. Eine zusätzliche Funktionalität wird durch die veränderte Farbe des Lichtstrahls herbeigeführt. Hierbei kann nur der Strahl reflektiert werden, wie auch der Spiegel aufweist.
    
    \TODO{TODO: Aussehen des Spiegels}
    
    \TODO{TODO: verschiedene Farben $\rightarrow$ Spiegel Rot $=$ roter Strahl wird reflektiert}
    
    \item \textbf{Beweglicher Spiegel} \hfill \\
    Der bewegliche Spiegel ist ein Game Objekt, der auf einer definierten Linie verschoben werden kann. Der Spiegel selbst kann allerdings nicht rotiert werden. Durch anklicken des Game Objektes wird der Spiegel ausgewählt, der dann wiederum durch wischen bewegt werden kann.
    
    \TODO{TODO: Aussehen der Linie}
    
    \TODO{TODO: Aussehen Spiegel}
    
    \item \textbf{Fester Spiegel} \hfill \\
    Der feste Spiegel ist ein unbewegliches Game Object, welches lediglich den Lichtstrahl reflektiert. Dieser Spiegel stellt keine weiteren Interaktionsmöglichkeiten zur Verfügung.
    
    \TODO{TODO: Aussehen des Spiegels}
\end{enumerate}

\item \textbf{Prisma} \hfill \\
Das Prisma ist ein Game Objekt, das den Lichtstrahl in seine Spektralfarben aufteilt. Hierbei ist der Auftreffpunkt ausschlaggebend, da dieser den Verlauf der weiteren Lichtstrahlen definiert. Eine direkte Interaktion über Berührung ist nicht möglich.

Das Prisma hat die Form eines Edelsteins, welcher bunt schimmert. Die schimmernden Farben setzen sich aus den Spektralfarben zusammen, die sich langsam abwechseln.

\item \textbf{Wasserfall} \hfill \\
Ein Wasserfall weist eine ähnliche Funktionalität wie ein Prisma auf. Sobald der Lichtstrahl die Wasseroberfläche (eigentlich den Nebel vom Wasserfall) trifft wird der Strahl in seine Spektren aufgeteilt. Diese neu entstandenen Lichtstrahlen (insgesamt sieben: Rot, Orange, Gelb, Grün, Hellblau, Indigo und Violett) fallen in einem Bogen zum Erdboden. Es ist egal in welchem Winkel der Strahl auf den Wasserfall trifft, genau an der Schnittstelle beider Elemente wird der Regenbogen generiert. Der Wasserfall selbst besitzt keine weiteren Interaktionsmöglichkeiten.

\TODO{TODO: Beschreibung Aussehen des Wasserfalls}

\TODO{TODO: Wasserfall (und ggf. Nebel) wird animiert}

\item \textbf{Türmechanismus} \hfill \\
Es existieren Tore, die den Weg des Lichtstrahls versperren. Diese Tore müssen über einen Schalter geöffnet werden. Der Türmechanismus unterscheidet sich in zwei Typen:
\begin{enumerate}
\item \textbf{Einmaliges Anleuchten des Schalters}\hfill \\
Hierbei muss der Lichtstrahl nur eine kurze Zeit auf den Schalter gerichtet sein. Dieser lädt sich innerhalb von zwei Sekunden auf, danach wird das Tor durch eine Animation dauerhaft geöffnet.
\item \textbf{Dauerhaftes Anleuchten des Schalters}\hfill \\
Hierbei muss der Lichtstrahl dauerhaft auf den Schalter leuchten, sodass sich das Tor öffnet und auch geöffnet bleibt. Der Schalter lädt sich eine kurze Zeit auf. Dann erfolgt eine Animation, wie sich das Tor öffnet, bzw. schließt.
\end{enumerate}
Beiden Typen können nicht miteinander kombiniert werden (für eine Tür), da beide Typen das selbe Design aufweisen. Es können allerdings mehrere Schalter eines Typs für ein Tor verwendet werden.

\TODO{TODO: Beschreibung Aussehen der Schalter $\rightarrow$ Sozalzelle}

\TODO{TODO: Beschreibung Aussehen der Tore $\rightarrow$ Textur aus Stein/Holz/Blätter-platten, abhängig vom Level}

\TODO{TODO: Beschreibung Aussehen der Türöffnungsanimation}

\item \textbf{Wasser} \hfill \\
Das Medium Wasser ist zuständig für eine weitere physikalische Besonderheit: die Brechung des Lichts. Hierbei soll der Lichtstrahl am Auftrittspunkt M im Wasser gebrochen werden. Zeitgleich erfolgt eine Reflexion auf der Oberfläche. Man kann den Strahl nur in, bzw. aus dem Wasser führen, wobei immer die physikalischen Gegebenheiten beachtet werden. Eine direkte Interaktion ist nicht vorgesehen.

Das Wasser ist nur aus der Seitenperspektive sichtbar. Hierbei wird eine waagerechte Linie durch das Bild gezogen und das Wasser etwas bläulich eingefärbt. Dadurch ist es als solches erkennbar.

\item \textbf{Glas} \hfill \\
Glas weist das selbe Prinzip wie Wasser auf. Der Lichtstrahl trifft auf Glas, woraufhin der Strahl gebrochen und reflektiert wird.

\TODO{TODO: Interaktionsmöglichkeiten? / braucht man Glas?}

\TODO{TODO: Aussehen des Glases beschreiben}

\item \textbf{Hintergrund} \hfill \\
Der Hintergrund dient der Umgebungsgestaltung. Hierbei sind also keinerlei Interaktionsmöglichkeiten verbunden. 

Der Hintergrund besteht aus Felsen, Berge, Bäume und anderen Pflanzen. Durch geschicktes anordnen der Assets aus dem 2- und 3D-Bereich entsteht der Eindruck einer 3D-Umgebung. Zudem wird der Hintergrund in seinen natürlichen Farben dargestellt.

\item \textbf{Vordergrund} \hfill \\
Auch der Vordergrund dient der Umgebungsgestaltung. Hierbei sind allerdings auch Interaktionsmöglichkeiten gegeben. Einige dieser sind die Steuerung des Lichtstrahl sowie der anderen Game Objekte.

Assets wie Felsen, Gräser, Steine und andere dienen als Wände und somit als Begrenzung des Spielfeldes. Die Assets befinden sind in einem minimalistischen Zustand, das heißt, sie sind nur als schwarze und graue Flächen zu erkennen. Erst bei einsammeln der Collectables wird der Vordergrund in Farbe gehüllt.

\item \textbf{Lichtstrahl} \hfill \\
Der Lichtstrahl ist das eigentliche Spielelement, wobei eine direkte Interaktion nicht möglich ist. Das wird mit der linearen, gleichmäßigen Ausbreitung des Lichts begründet. Die Steuerung der Lichtstrahl erfolgt vom jeweiligen Ausgangspunkt. Das bedeutet, das Game Objekt, welches den Lichtstrahl erzeugt, bzw. reflektiert sind die eigentlichen Steuerelemente. Der Strahl selbst ist immer geradlinig und wird ggf. an Spiegel reflektiert oder gebrochen oder endet an den Wänden. Trifft er auf Wände oder auf nicht reflektierende Game Objekte wird eine Animation erzeugt, die das Strahlende zeigt. Er wird unterschieden in zwei Typen:
\begin{enumerate}
\item \textbf{Normaler Lichtstrahl} \hfill \\
Der Normale Lichtstrahl wird als weißer Strahl mit einer Animation dargestellt. Diese Animation lässt den Strahl "{}lebendig"{} wirken. Zudem werden Partikeleffekte verwendet, um den Strahl optisch aufzuwerten.
\item \textbf{Farbiger Lichtstrahl} \hfill \\
Der farbige Lichtstrahl besitzt genau die selben Eigenschaften wie der normale. Der einzige Unterschied besteht darin, das der Strahl nicht weiß, sondern eine andere Farbe besitzt. Er kann folgende Farben annehmen: \textcolor{red}{\textbf{Rot}}, \textcolor{orange}{\textbf{Orange}},  \textcolor{yellow}{\textbf{Gelb}}, \textcolor{green}{\textbf{Grün}}, \textcolor{cyan}{\textbf{Hellblau}}, \textcolor{indigo}{\textbf{Indigo}} und \textcolor{violet}{\textbf{Violett}}
\end{enumerate}

\item \textbf{Startpunkt} \hfill \\
Der Startpunkt eines jeden Level erzeugt den Lichtstrahl. Durch einmaligen berühren des Startpunktes aktiviert sich dieser mit einer Animation und der Lichtstrahl wird dauerhaft erzeugt. Ein erneutes auswählen lässt den Lichtstrahl nicht verschwinden, sondern wählt das Game Objekt aus, sodass der Lichtstrahl um 360\degree rotiert werden kann. Der Strahl steht somit im direkten Zusammenhang zum Startpunkt. Die Rotation wird durch wischen oder tippen auf eine beliebige Stelle im Spiel erzeugt. Der Lichtstrahl und somit die Ausrichtung des Startpunktes orientiert sich an der Position an dem der Bildschirm berührt wird.

Der Startpunkt soll in Form einer Sonnenblume dargestellt werden. Diese befindet sich zuerst in einem nicht-erblühten Zustand. Beim Antippen des Game Objektes startet eine Animation, die nicht länger als zwei Sekunden dauert. Lichtpartikel sammeln sich aus der Luft an der Sonnenblume, zusätzlich erblüht diese. Die Lichtpartikel verschwinden und ein Lichtstrahl generiert sich aus der Blüte. Die Blüte kann um 360\degree rotiert werden.

\item \textbf{Zielpunkt} \hfill \\
Der Zielpunkt stellt das Ziel eines jeden Levels dar. 
Sobald das Ziel mit dem Licht angestrahlt wird, ist das Levels gelöst. Durch das Berühren des Zielpunkts kann der Spieler das Level beenden. Falls dieser alle Collectables eingesammelt hat erstrahlt auch der letzte Levelabschnitt in Farbe und der Endscreen dieses Levels wird angezeigt.

Je nach Fortschritt und Zustand des Baums wird der Zielpunkt anders dargestellt. Dies beginnt mit einem Samen und verändert sich im laufe des Spiels zu einem Sprössling und weiter zu einem Baum, welcher das Licht benötigt um weiter zu wachsen.

\item \textbf{Checkpoint} \hfill \\
Der Checkpoint innerhalb eines Levels verbindet die einzelnen Levelabschnitte miteinander und erlaubt das weiter Scrollen des Bildschirms zum nächsten Teil des Levels.
Sobald der Lichtstrahl auf den Checkpoint trifft leuchtet der Checkpoint auf und ermöglicht dem Spieler durch eine Berührung dieses Elements in den nächsten Abschnitt zu gelangen. Wenn alle Collectables innerhalb des Abschnitt gesammelt wurden, erstrahlt dieser in Farben. 

Der Kontrollpunkt soll in Form einer kleineren Sonnenblume dargestellt werden. Diese befindet sich zuerst in einem nicht-erblühten Zustand. Die Pflanze erblüht beim Berühren während der Lichtstrahl darauf leuchtet. 

\item \textbf{Collectable-Pflanze} \hfill \\
Die Col.-Pflanze wird als Vorder- und Hintergrund Stilmittel eingesetzt, um die Umgebung zu gestalten. Dieses Game Objekt steht immer im direkten Zusammenhang mit einem Collectable. Eine Interaktion ist allerdings nicht möglich.

Die Col.-Pflanze ist eine Blume / Blüte, die zu Beginn eine Knospe ist. Beim "{}Einsammeln"{} eines Collectables fällt dieser auf die Col.-Pflanze, welche wiederum aufblüht und die Umgebung in Farbe hüllt. 

\TODO{TODO: Überlegen welche Pflanze schön dazu passt.}

\item \textbf{Wand} \hfill \\
Die Wände dienen als Spielfeldbegrenzung oder als Hindernisse, welche den Lichtstrahl blockieren. Mit den Wänden ist keine Interaktion möglich. Es ist dem Lichtstrahl nicht möglich durch die Wand zu strahlen.

Wände können durch unterschiedlichste Assets zusammengesetzt sein. Diese sind unter anderen Felsen, Steine und Bäume. Auch die Wände befinden sich in einem minimalistischen Zustand und werden daher in Schwarz / Grau dargestellt. Eine farbige Gestaltung erfolgt bei einsammeln eines Collectables.

\item \textbf{Collectables} \hfill \\
Collectables sind Game Objekte, welche gesammelt werden können. Das erfolgt allerdings nicht durch Antippen, sondern durch anstrahlen des Lichts. Sie werden erst bei Abschluss des Levels oder des Bereiches eingesammelt, unter der Bedingung, dass noch immer das Collectable angeleuchtet wird. Mit Ihnen ist keine direkte Interaktion möglich. 

Das Collectable hat die Form eines Wassertropfens. Dieser ist halbdurchsichtig und bläulich eingefärbt. 

%\item \textbf{Weltenbaum} \TODO{TODO} \hfill \\
%
%Der Weltenbaum dient als Levelauswahlmenü und wird je nach Anzahl der freigeschalteten Level größer und umfangreicher gestaltet. Folgende Stadien des Baumes sind möglich Samen/Knospe/Strauch/Baum etc. Das Ziel des Spiels ist es den Baum zum wachsen zu bringen und ihn in seine vollste Pracht zu bringen.

\end{enumerate}
\subsection{Grafikdesign}
\subsubsection{Grafikstil des Spiels}
Das Spiel ist in einem schlichten, minimalistischen Design gehalten. Zu Beginn sind nur wenige Farben erkennbar, diese sind unter anderem schwarz-weiß, verschiedene Graustufen und Grüntöne. Durch Einsammeln der Collectables und Aktivierung des Checkpoints oder Zielpunktes werden in einem definierten Bereich alle Assets farblich gestaltet. Sind alle Collectables in einem Level eingesammelt werden auch alle Assets im gesamten Bereich farblich dargestellt.

\subsubsection{Verwendete Assets}
Für die Spielelemente sowie das Design der Level werden u.a. die Assets "{}Stylized Jungle Pack"{}\cite{jungle} und {}"2D Forest Pack"{} \cite{forest} vom schwedischen Grafikdesigner Mikael Gustafsson, welche über den Assetstore von Unity erworben wurden, verwendet. Diese entsprechen unseren Vorstellungen über das Spiel und tragen viel zur Immersion der Spielwelt bei.

\subsection{Musik}
\subsection{Spezifikation}
\begin{description}
\item [Horizontale Verwendung] \hfill \\
Das Spiel wird nur in der Horizontalen Perspektive des Geräts dargestellt.

\item [Automatische Speicherung] \hfill \\
Das Spiel soll nach jedem Level automatisch gespeichert werden. Dafür wird lokal eine Speicherdatei erzeugt, die im Game Ordner versteckt angelegt wird. Die bereits vorhandene Optional wäre den Spielstand in die Googleplay Cloud zu laden.
\end{description}

\subsection{Spielmechaniken}

\newpage
\section{Implementierung}

\newpage
\section{Interne Qualitätssicherung}

\newpage
\section{Testphase (Beta)}

\newpage
\section{Optimierung des Game Plays}

\newpage
\section{Release Version 1.0 im Google Play Store / IOS}

\newpage

\bibliographystyle{unsrt}
\bibliography{literatur}
\addcontentsline{toc}{section}{Literaturverzeichnis}



\section{Linkliste}

Unity / Engine / Programmierung
\begin{itemize}
\item https://forums.unrealengine.com/showthread.php?75125-Unreal-Vs-Unity-5-for-mobile-2d-games-is-unreal-to-powerful
\item https://create3dgames.wordpress.com/2015/09/07/unity-5-vs-unreal-engine-4/
\item http://answers.unity3d.com/questions/56140/how-to-make-a-main-menu.html
\item http://web.fhnw.ch/technik/projekte/computervision/fhbb/studierendenprojekte/2004/pa-gpu-programmierung/licht-shader
\item http://unity-buch.de/screencasts/videos/tutorials
\item http://code.tutsplus.com/articles/effectively-organize-your-games-development-with-a-game-design-document--active-10140
\item http://www.gdv.informatik.uni-frankfurt.de/diplomarbeiten/pdf/Diplomarbeiten\_final/Diplomarbeit\_mraab.pdf
\item https://docs.unrealengine.com/latest/INT/Engine/Physics/index.html
\item http://box2d.org/about/
\item http://docs.unity3d.com/Manual/PhysicsSection.html
\item https://docs.unrealengine.com/latest/INT/Engine/Paper2D/index.html
\item https://docs.unrealengine.com/latest/INT/Engine/Rendering/index.html
\item http://wiki.blender.org/index.php/Doc:DE/2.6/Manual/Game\_Engine
\item https://www.blender.org/manual/physics/index.html
\item https://www.blender.org/manual/game\_engine/introduction.html
\item https://www.blender.org/support/faq/
\item http://www.blender.org/about/license/
\item http://wiki.blender.org/index.php/Doc:DE/2.6/Manual/Physics
\item http://docs.unity3d.com/Manual/PhysicsSection.html
\item http://docs.unity3d.com/Manual/Graphics.html
\item http://unity3d.com/get-unity
\item https://www.unrealengine.com/what-is-unreal-engine-4
\item http://www.vis.uni-stuttgart.de/plain/seminare/computerspiele/zhenpeng/Ausarbeitung.pdf

\end{itemize}

Studien
\begin{itemize}
\item \circledmark\quad http://www.gulli.com/news/17508-uni-wuerzburg-computerspiele-foerdern-nachweislich-die-intelligenz-von-kindern-2011-11-14
\item \circledmark\quad http://www.uni-wuerzburg.de/sonstiges/meldungen/single/artikel/intelligen-3/
\item \circledmark\quad http://www.bpb.de/gesellschaft/medien/computerspiele/63725/probleme-loesen
\item \circledmark\quad http://www.fr-online.de/wissenschaft/lernspiele-spielend-physik-lernen,1472788,11542622.html
\item http://www.physik.org/stichwort.html
\item \circledmark\quad https://www.medienpass.nrw.de/de/inhalt/was-lernt-man-beim-computerspielen
\item https://www.fbi.h-da.de/fileadmin/personal/b.harriehausen/Forschungsprojekte/SigmaStar/art/Querschnitt\_20.pdf
\item \circledmark\quad http://www.psychologie-aktuell.com/fileadmin/download/esp/2-2011\_20110927/lenhard.pdf
\item \circledmark\quad http://www.psychometrica.de/denkspiele\_evaluation.html
\item \circledmark\quad http://www.aappublications.org/content/early/2015/04/26/aapnews.20150426-4?sso=1\&sso\_redirect\_count=1\&nfstatus=401\&nftoken=00000000-0000-0000-0000-000000000000\&nfstatusdescription=ERROR\%3a+No+local+token
\item \circledmark\quad http://pediatrics.aappublications.org/content/134/5/849
\item http://truben.no/latex/bibtex/\#
\end{itemize}


Assets
\begin{itemize}
\item 2D Jungle Pack (https://www.assetstore.unity3d.com/en/\#!/content/43458)
\item 2D Forest Pack (https://www.assetstore.unity3d.com/en/\#!/content/45564)
\item Hand Painted Jungle (https://www.assetstore.unity3d.com/en/\#!/content/43609)
\item Stylized Jungle Pack (https://www.assetstore.unity3d.com/en/\#!/content/40372)
\item Stylized Nature Pack 
https://www.assetstore.unity3d.com/en/\#!/content/37457
\item 2D Forest Pack
https://www.assetstore.unity3d.com/en/\#!/content/41171
\end{itemize}

\newpage
\section{Game Design}

\subsection{Ziel des Spiels}
Das Ziel des Spiels ist es, einen Samen zu einem Weltenbaum zu entwickeln. 

\subsection{Erreichen des Ziels}
In einem Puzzle-Spiel müssen mehrere Level erfolgreich durchlaufen und Collectables gesammelt werden. Das erfolgt durch Steuerung eines Lichtstrahls, der auf das Ziel-Objekt des Levels treffen muss. Die Freischaltung der Level erfolgt durch Abschließen des vorherigen Levels und durch Besitzen einer definierten Anzahl Collectables. Zur vollständigen Entwicklung des Baumes werden 50 Level durchlaufen.

\subsection{Level Screen / Bereiche}
Das Konzept dieses Spiels beruht auf die Verwendung mehrerer Screens. Hierfür wird es einen Tutorialbereich geben, der die grundlegende Mechanik des Spiels erläutert. Dieses Tutorial und die ersten Level wird je 1 Screen bereit gestellt. In den weiterführenden Level steigt die Anzahl der Screens, die zum erreichen des Zielpunktes notwendig sind. Mehrere Screens zeichnen sich ähnlich wie ein Jump'n Run-Game aus, da sich die Kamera bzw. das Level automatisch bewegt. Bei diesem Spiel ist die Kamera und auch der Bereich fixiert. Erst beim erreichen des Lichtstrahls auf einem Checkpoint und dessen Aktivierung, verschiebt sich die Kamera zum nächsten Bereich des Levels. Eine Aufsplittung der Level und wie viele Bereiche diese haben sind in Tabelle \ref{screenLevel} aufgeführt.

\begin{table}[H]
\centering
\begin{tabular}{c|c}
Levelbereich & Anzahl der Screens / Bereiche \\ \hline
1 - 10  & 1 Screen \\ 
11 - 30 & 2 Screens \\
31 - 50 & 3 Screens      
\end{tabular}
\caption{Anzahl Screens in Levelbereiche}
\label{screenLevel}
\end{table}

%%%%%%%%%%%%%%%%%%%%%%%%%%%%% BEGIN OLD
Hierbei werden mehrere Level durchlaufen. Neue Level werden durch den Abschluss des vorherigen Levels und  eine definierte Anzahl an Collectables freigeschaltet. Im Durchschnitt sind 3 Coll. gedacht. Bei mehreren Screens können die Coll. variieren. Insgesamt sind 50 Lvl für ein Weltenbaum angedacht. Es werden maximal 3 Screens pro Lvl verwendet.
Collectables werden voraussichtlich als Wassertropfen dargestellt. Bei abschließen des Lvls / Bereichs fallen die Wassertropfen auf Pflanzen, die erblühen und ggf. die Map erhellen.

Collectables notwendig für nächstes Lv
\begin{description}
\item [Lv 1 - 5] 0\%, 3C je Lv 
\item [Lv 6 - 10] 30\%, 3C je Lv
\item [Lv 11 - 30] 50\%
\item [Lv 31 - 45] 70\%
\item [Lv 46 - 50] 70\%   (von der bisherigen Gesamtanzahl)
\end{description}
%%%%%%%%%%%%%%%%%%%%%%%%%%%%%%% END OLD

\subsection{Design des Spiels}
Das Spiel ist in einem schlichten, minimalistischen Design gehalten. Zu Beginn sind nur wenige Farben erkennbar, diese sind unter anderem schwarz, weiß, verschiedene Graustufen und Grüntöne. Durch Einsammeln der Collectables und Aktivierung des Checkpoint / Zielpunktes werden in einem definierten Bereich alle Assets farblich gestaltet. Dabei nehmen die Assets ihre ursprüngliche Farbe an, welche sie von Beginn an besaßen. Werden alle Collectables in einem Bereich eingesammelt werden auch alle Assets im gesamten Bereich farblich dargestellt.

Zudem ist der Stil des Spiels im Comic-Style gehalten. Die Aufbereitung der einzelnen Elemente (Wände, Pflanzen etc.) können allerdings alle möglichen Asset-Arten beinhalten. Die Kombination der unterschiedlichen Assets erzeugen den gewünschten Comic-Look.

\subsection{Wie wird der Checkpoint aktiviert / verwendet?}
Der Lichtstrahl muss den Ch-P berühren. Dazu muss der User den Point anklicken um in den nächsten Screen zu gelangen. Dadurch wird ein zufälliges aktivieren des Pointes verhindert.



\subsection{Spielelemente}
\begin{description}

\item [$\circledmark\quad$Spiegel] \hfill \\
Ein Spiegel ist ein Game Objekt, welches den Lichtstrahl reflektiert. Hierbei wird die korrekte Reflexion nach physikalischen Gesetzen verwendet. 

\begin{description}
    \item [$\circledmark\quad$rotierender Spiegel] \hfill \\
    Ein Spiegel welcher in einem Winkel von 140 Grad rotiert werden kann (Entgegengesetzt zur Wand).
    
    \item [$\circledmark\quad$Spiegel mit Farbeinschränkungen] \hfill \\
    Ein Spiegel, welcher nur bestimmte Lichter reflektiert.
    
    \item [$\circledmark\quad$beweglicher Spiegel] \hfill \\
    Ein Spiegel welcher auf einer Horizontalen/Vertikalen Bahn durch wischen bewegt werden kann.
    
    \item [$\circledmark\quad$feste Spiegel] \hfill \\
    Spiegel die keine Interaktionsmöglichkeiten besitzen.
    
\end{description}

\item [$\circledmark\quad$Prisma] \hfill \\
Wird als Kristall/Edelstein welcher Bunt schimmert(Farbe wechselt) dargestellt. Bei eintreffen des Lichtstrahls wird das Licht gebrochen und in die einzelnen Spektralfarben aufgeteilt. Das Prisma ist ein festes Objekt mit welchem keine weiteren Interaktionsmöglichkeiten bestehen.

\item [$\circledmark\quad$Wasserfall] \hfill \\
Wasserfälle funktionieren ähnlich wie Prismen, diese teilen das Licht durch das Wasser  in einen Regenbogen. Der Lichtstrahl muss im richtigen Winkel im Wasser auftreten um durch den Bogen die Checkpoints oder Türmechanismen zu aktivieren. Die Wasserfälle sollten animiert sein.

\item [$\circledmark\quad$Türmechanismus] \hfill \\
Türschalter werden als kleine Solarzellen dargestellt, welche Energie aufladen und nach kürzester Zeit ein Aufzugstor öffnen. Diese können aus einem oder zwei Mechanismen bestehen um ein Tor zu öffnen. Die Tore bestehen aus Stein/Holz/Blätter-platten die Textur ist Abhängigkeit der Level auszuwählen.

\item [Flüssigkeiten] (Optional) \hfill \\
Hierbei wird speziell auf das Wasser eingegangen. Dieses soll in der Seitenperspektive dargestellt werden. (Wahrscheinlich nur in der Höhle anwendbar) Der auf das Wasser treffende Lichtstrahl wird entsprechend physikalischer Bedingungen gebrochen. Beim austreten des Strahl aus dem Wasser werden ebenfalls physikalische Begebenheiten beachtet.

\item [$\circledmark\quad$Hintergründe] \hfill \\
Hintergründe dienen der Umgebungsgestaltung und den Eindruck einer 3D Umgebung. Diese bestehen aus Wänden, Felsen, Bergen, Bäumen, Pflanzen, Gräsern und Partikeleffekte. Der Hintergrund wird in normalen Farbe dargestellt.

\item [$\circledmark\quad$Vordergrund] \hfill \\
Vordergründe dienen der Umgebungsgestaltung und den Eindruck einer 3D Umgebung. Dieser besteht größtenteils aus Wänden , Pflanzen, Steinen, Felsen und Gräsern. Die Darstellung ist bis zum Aufdecken der Karte in minimalistischen Farben zu halten. (Farben mit wenig Intensität). Nach einsammeln aller Collectables wird die Karte als aufgedeckt behandelt. 

\item [$\circledmark\quad$Lichtstrahl] \hfill \\
Nachdem auswählen der jeweiligen Lichtquelle/Reflexionspunkt lässt der Lichtstrahl durch wischen auf dem Bildschirm in eine Richtung lenken.
Strahlt bis Ende des Bildschirms oder bis er auf ein Objekt trifft
Reflektiert beim Eintreffen auf ein Spiegelähnliches Objekt
Bricht in verschiedene Farbstrahlen beim Eintreffen auf ein Prisma.
Beim durchdringen von Wasserfällen verhält sich der Lichtstrahl  wie beim Prisma. und eine Animations für ein Regenbogen wird ausgeführt
Bei Blenden von Lebewesen werden Animation durchgeführt
Lichtstrahl beinhaltet Partikeleffekte für die optische Aufwertung
Der Lichtstrahl ist weiß

\item [Startpunkt] \hfill \\
Der Startpunkt eines jeden Levels ist eine Sonnenblume. Diese ist zunächst als einfache Pflanze zu sehen und wird durch natürliche Sonnenstrahlen beleuchtet. Durch anklicken der Pflanze wird eine Animation eingeleitet, welche die Sonnenblume erblühen lässt. Die Animation soll nicht länger als 1-2 Sekunden dauern, da sie kein direkten Bestandteil des Spielvergnügens darstellt. Ist sie vollends aufgeblüht erzeugt der Blütenkopf von selbst aus ein Lichtstrahl. Dieser Lichtstrahl dient als Spielmedium. 

Die Steuerung des Lichtstrahls erfolgt durch einmaliges antippen des Blütenkopfes. Dadurch lässt sich der Kopf um 360 Grad drehen, entsprechend dazu verändert sich ebenfalls die Richtung des Lichtstrahls.

\item [$\circledmark\quad$Zielpunkt] \hfill \\
Der Zielpunkt ist der Baum  selbst je nach Zustand des Baums/Fortschritt des Spiels (Samen/Sprössling/Baum). 
Sobald der Zielpunkt angestrahlt wird und das Rätsel gelöst ist kann durch ein Klick auf den Zielpunkt das Level beendet werden. Danach sieht man den Endscreen des Levels.
\item [Checkpoint]\hfill \\

\item [$\circledmark\quad$Pflanzen] \hfill \\
Pflanzen werden als Vordergrund/Hintergrund eingesetzt um die Umgebung zu gestalten. Diese werden in Schwarzweißer Optik und möglichst mit leichten Animationen dargestellt. Falls alle Collectables eingesammelt wurden, werden diese im entsprechenden Bereich nun in Farbe angezeigt.

\item [$\circledmark\quad$Wände] \hfill \\
Die Wände des Spiels werden durch unterschiedliche Assets begrenzt. Hierbei dienen Felsen, Bäume, Sträucher, Gräser, Lianen, Steine, Holz, Pflanzen. Sobald der Lichtstrahl auf eines dieser Elemente trifft, wird eine Animation für das Lichtende abgespielt. Der Lichtstrahl durchdringt diese Elemente in der Regel nicht, es sei denn sie befinden sich im Hintergrund / Vordergrund.

\item [Felsen]\hfill \\

\item [$\circledmark\quad$Hügel] \hfill \\
Berge und Hügel dienen zur Gestaltung des Hintergrunds.

\item [$\circledmark\quad$Collectables] \hfill \\
Hierbei handelt es sich um Wassertropfen. Diese können eingesammelt werden, indem der Lichtstrahl über diese gesteuert werden. Damit die Coll. eingesammelt werden, muss der Lichtstrahl zum Abschluss des Lv/Bereichs noch immer auf die Coll. verweilen, deshalb müssen diese so platziert werden dass auch bei einem Durchgang alle eingesammelt werden können.

\item [$\circledmark\quad$Animation] \hfill \\
Animationen werden eingesetzt um das Feedback einer Interaktion des Spieler zu visualisieren und den optischen Eindruck der Umgebung zu verbessern.

\item [$\circledmark\quad$Tierchen] \hfill \\
Tierchen dienen als Funfact und sorgen für eine authentische Umgebung. Optional werden Eastereggs durch kleine Animationen der Tierchen eingebaut falls diese vom Lichstrahl erwischt werden. ( Maulwurf holt Sonnenbrille, Spinne verkriecht sich)

\item [$\circledmark\quad$Weltenbaum] \hfill \\
Der Weltenbaum dient als Levelauswahlmenü und wird je nach Anzahl der freigeschalteten Level größer und umfangreicher gestaltet. Folgende Stadien des Baumes sind möglich Samen/Knospe/Strauch/Baum etc. Das Ziel des Spiels ist es den Baum zum wachsen zu bringen und ihn in seine vollste Pracht zu bringen.

\item [Hauptmenü] \hfill \\
siehe Zeichnung

\item [Optionsmenü] \hfill \\
siehe Zeichnung (Play und Optionen Buttons)

\item [Ansicht der Collectables] \hfill \\
Siehe Zeichnung 

\item [Abschlussdarstellung] \hfill \\
Siehe Zeichnung
Es ist möglich das Level Neuzustarten, das nächste Level zu starten, in die Levelansicht zu wechseln und/oder das gerade gespielte Level komplett zu betrachten. 

\item [Horizontale Verwendung] \hfill \\
Das Spiel wird nur in der Horizontalen Perspektive des Geräts dargestellt.

\item [Automatische Speicherung] \hfill \\
Das Spiel soll nach jedem Level automatisch gespeichert werden. Dafür wird lokal eine Speicherdatei erzeugt, die im Game Ordner versteckt angelegt wird. Die bereits vorhandene Optional wäre den Spielstand in die Googleplay Cloud zu laden.

\end{description}


\subsection{Setting}
\begin{description}
\item [Höhle]
\item [Erdoberfläche/Wiese]
\item [Wald]
\item [Baumkrone]
\end{description}

\subsection{optional}
\begin{description}
\item [Bauminneres]
\item [Himmel]
\end{description}
\newpage



%% Der Kreis mit dem Haken in der Mitte
%%  $\circledmark\quad$
\section{TODOs}

\subsection{Zwischenstand}
\begin{itemize}
\item Ausarbeitung Prototyp Level (vllt erste Konzepte für die ersten LV)
\item Ausarbeitung Assets (Teil 1)
\item Game Design Dokument
\item Komponenten anhand Gamedesign und Gamification Richtlinien erläutern (Genre, Direktes Feedback, Belohnungssysteme usw.)
\item Ausarbeitung Analyse ob Spiele lehren + Text verfassen
\item Ausarbeitung Analyse der versch. Game Engines (Vor-/ Nachteile) + Wahl welcher Engine + Text verfassen

\item Einarbeitung in Unity + Sichtung der notwendigen Physik skripte
\item Ermittlung der benötigten Assets (Grundlage)
\item Ermittlung der benötigten Effekte (wie Partikel)
\item Animationen definieren
\item Ermittlung benötigter physikalischer Formeln (Gegebenheiten)
\item Einarbeitung dokumentieren



\item Abgabe Anfang - Mitte Januar
\item 80\% des wissenschaftlichen Bereiches sollte abgehakt sein
\item auf jeden Fall God of Light mit aufnehmen als Inspiration
\item Analyse bzw. auch das Auseinandersetzen mit Physik und Game Engine anfangen Überarbeiten
\item $\rightarrow$ da kann man auch noch wissenschaftlichen Kran rausholen aus der Doku
\item ich (Matze) schätzt das wir ca 20-30\% wissenschaftlichen Teil erst erledigt haben
\item
\end{itemize}

\subsection{15.10.2015}
\begin{itemize}
\item \circledmark\quad  Studien bzgl Spielen (Text verfassen) $\rightarrow$ Matze
\item \circledmark\quad  Engines vergleichen (Engines verfassen)$\rightarrow$ Roman
\item \circledmark\quad  Spieldesign (Gutes Design/NoGos) $\rightarrow$ beide
\item \circledmark\quad  Assets in Unity $\rightarrow$ Matze
\item \circledmark\quad Fragen wegen Budget $\rightarrow$ Matze $\rightarrow$ Muss individuell beim Freudenmann geprüft werden
\end{itemize}

\subsection{22.10.2015}
\begin{itemize}
\item \circledmark\quad Game mechanics besprechen (Teils) $\rightarrow$ zusammen
\item \circledmark\quad Bücher weiter lesen und ggf. Passagen rausschreiben $\rightarrow$ beide
\item \circledmark\quad nebenher mit Unity befassen $\rightarrow$ beide
\item \circledmark\quad Game Elemente besprechen und aufschreiben (Teil 1) $\rightarrow$ zusammen
\end{itemize}

\subsection{29.10.2015}
\begin{itemize}
\item \circledmark\quad  Zeitplan erstellen (vorerst bis Ende Dezember)
\item \circledmark\quad  Bücher auslesen + rausschreiben $\rightarrow$ beide
\item Unity mehr einarbeiten
\end{itemize}

\subsection{05.11.2015}
\begin{itemize}
\item \circledmark\quad Mögliche Gliederungspunkte / Überschriften für Dokument bestimmen $\rightarrow$ beide
\item siehe vom 29.10. TODOs müssen noch durchgeführt werden
\end{itemize}

\subsection{12.11.2015}
\begin{itemize}
\item entfallen
\item Definition der Game-Elemente + Aussehen, Funktion, Interaktion $\rightarrow$ beide

\end{itemize}

\subsection{19.11.2015}
\begin{itemize}
\item Game Mechanics zu den Game Elementen genauer definieren
\item \circledmark\quad  1. Version Analysetext zu Spiele und Lernen $\rightarrow$ Matze
\item \circledmark\quad  1. Version Game Engines (Vor- Nachteile) $\rightarrow$ Roman
\item Game Design (Dokument) überarbeiten auf Vollständigkeit und Logik etc.
\end{itemize}

\subsection{26.11.2015}
\begin{itemize}
\item Assets Liste auswählen
\item FB Text für Assets schreiben
\item \circledmark\quad Texte in Latex konvertieren 
\item \circledmark\quad  Meilenstein Sonntag $\rightarrow$  unsere Texte \bfseries{müssen} fertig werden!!!!!  
\item \circledmark\quad  Emily kontaktieren
\end{itemize}


\subsection{03.12.2015}
\begin{itemize}
\item Roman fragt nach entgültigen Abgabetermin Zwischenstand
\item Mikael Gustafsson (asset publisher mikaelgustaf@gmail.com) anschreiben wegen kostenloser assets
\item Grafiker gesucht text in FB und Foren nachfragen
\item 
\end{itemize}

\subsection{Optimierungen}
\begin{itemize}
\item Gamesdesign Dokument strukturieren in Optisch/Funktion(interaktiv) (Spielelemente)
\item Genres und Grafikstil erklären (warum diese Entscheidung)
\item 
\end{itemize}

\subsection{10.12.2015}
\begin{itemize}
\item 
\item 
\item 
\item 
\end{itemize}

Abgabe Zwischenstand: 12.12.2015

%%%%%%%%%%%%%%%%%%%%%%%%%%%%%%%%%%%%%%%%%%%%%%%
%% Wenn man ein TODO abgeschlossen hat        %
%% bitte \circledmark\quad dahinter schreiben %
%% \circledmark\quad                          %
%%%%%%%%%%%%%%%%%%%%%%%%%%%%%%%%%%%%%%%%%%%%%%%

%%%%%%%%%%%%%%%%%%%%%%%%%%%%%%%%%%%%%%%%
%% Rechtspfeil                         %
%% $\rightarrow$                       %
%%%%%%%%%%%%%%%%%%%%%%%%%%%%%%%%%%%%%%%%

%% Analyse ob spielen das Lernen fördert / ob man durch spielen lernt
%% Analyse ob 2D oder 3D (Vor- NAchteile)
%% Game genre und Physikbereiche
%% Analyse Game Engines

%% Spielkonzept


\end{document}